\scanpage{157}
\mspage{106}

donc \unsure{tout} revient, à partir de \inserted{l'exemple donné} $X\mathrel{\overset{f}{\unsure{-|\to}}}Y$, à construire une flèche

\[
\unsure{\pi|\Pi}:\deleted{\unsure{X}}Y\longrightarrow U'\tag*{\unsure{$(x)$|$(\alpha)$}}
\]

dans $\mathrm{Cat}$ qui satisfasse les conditions

a) $\pi$ est un iso. local \qquad b) $U'$ est \inserted{$\unsure{W|\mathcal W}$-}asphérique.

On prendra alors $p=\pi f$

Mais par la façon dont $X$ a été construit : partir d'un $Y$ 1-connexe et non $\infty$-connexe, \unsure{par} \deleted{\unsure{et}} ailleurs quelconque, on est ramené finalement à trouver $(x)$ satisfaisant a), b) et de plus

c) $Y$ est 1-connexe, mais non $\infty$-connexe

(NB on suppose maintenant $W=W_\infty$). Il revient au $=$ de trouver

% vertical rule at left
\noindent\makebox[0pt][r]{\rule[-1.8\baselineskip]{0.45pt}{2.6\baselineskip}\hspace{.65em}}%
$U'$ objet \uline{asph.} de $\mathrm{Cat}$, $F\in\mathrm{Ob}\,{U'}^\wedge$ tels que $U'_{/F}$ soit 1-connexe et non $\infty$-connexe.

On songe par exemple à prendre simplement une sous-catégorie ouverte $U'_0$ de $U'$, qui soit 1-connexe et non $\infty$-connexe. Partons d'une telle catégorie $U'_0$ -- par exemple une réalisation \inserted{homotopique} dans $\mathrm{Cat}$, de la sphère $S^2$. On l'immerge \inserted{comme ouvert} dans un $U'$ asphérique, en lui ajoutant un objet final. On a gagné !
% archive sheet number: 157
